\cleardoublepage

\chapter{AI in Science -- Tool Instead of Truth}
\label{anhangC:ki}
\markboth{Appendix C}{Appendix C}

\subsection*{Motivation}

This book arose from the desire to consistently develop an idea
that runs through the entire history of mathematics and physics:
the world is not merely present, but structured --
and mathematics is the most precise form
in which this structure can be recognized.
The aim of this book is therefore
to make mathematics understandable not merely as a collection of techniques,
but as an inner order of number, form, motion, chance, and nature.

In the process, a new tool was used
that is becoming increasingly influential in scientific and literary work:
\textbf{artificial intelligence}\index{artificial intelligence},
specifically the language model \textbf{ChatGPT}\index{ChatGPT} by OpenAI.

But how can AI be used meaningfully in a book project
that depends precisely on conceptual clarity,
mathematical accuracy,
and intellectual independence?
And how can its use be disclosed openly
without blurring the author's responsibility?
This appendix gives a transparent insight
into the creation process of this book
and argues for a responsible use of AI as a tool --
not as truth.

\subsection*{What AI Can Do -- and What It Cannot Do}

AI-supported language models such as ChatGPT are powerful aids
for writing and structuring.
They can:

\begin{itemize}
	\item help formulate first drafts,
	\item smooth complex material linguistically,
	\item provide impulses for thinking or suggest outlines,
	\item offer stylistic alternatives.
\end{itemize}

What they \textbf{cannot} do, however:

\begin{itemize}
	\item \textbf{understand}\index{understanding} in the scientific sense,
	\item \textbf{verify} whether a formula has been derived correctly,
	\item \textbf{fully grasp physical concepts},
	\item \textbf{critically classify or evaluate sources}.
\end{itemize}

Therefore the rule is:
AI can \emph{support}, but it \textbf{cannot and must not replace the scientific process of gaining knowledge}.
Anyone who works with AI must still think for themselves --
and critically check every result.

\subsection*{How This Book Was Created}

The contents of this book --
from its structure to the physical explanations
and the mathematical derivations --
were conceived, researched, and responsibly developed by the author.
ChatGPT was used supportively in the following areas:

\begin{itemize}
	\item in \textbf{formulating individual passages},
	for example introductions, summaries, or didactic sections,
	\item for the \textbf{stylistic review} of technical sections,
	\item for the \textbf{development of outlines} in early working phases,
	\item for reflection on \textbf{comprehensibility}\index{comprehensibility}
	and reader guidance.
\end{itemize}

The decisive point is:
\textbf{All content-related statements, formulas, and interpretations were checked, questioned, revised, or rejected by the author.}
No AI replaced the subject-matter development of the physical argumentation.
The responsibility for the content,
selection,
and evaluation of all statements remained with the author throughout.

\subsection*{Ethical Questions and Scientific Responsibility}

The use of AI in science raises legitimate questions:

\begin{itemize}
	\item How much may be generated automatically
	without diluting authorship\index{authorship}?
	\item How should potential errors\index{error} be handled?
	\item How transparently must the use of AI be disclosed?
\end{itemize}

The answer lies in a basic principle of scientific work:
\textbf{responsibility}.
Anyone who uses AI remains responsible for the result --
regardless of whether individual formulations were suggested by a model.

In this sense, AI is not an author, but a tool.
It can accelerate processes,
but it cannot replace what science essentially requires:
\textbf{critical thinking, careful checking, and methodical work}.

\subsection*{Recommendations for Use in Research}

For researchers, teachers, and students,
this leads to a constructive path:

\begin{itemize}
	\item Use AI \textbf{deliberately and selectively} --
	for linguistic support, not for argumentation or proof.
	\item \textbf{Check every statement yourself} --
	especially in complex subject areas.
	\item \textbf{Disclose its use openly} when relevant --
	for example in prefaces, appendices, or submission declarations.
	\item Do not use AI for \textbf{deception}\index{deception}
	or as a fig leaf,
	but as support for presenting your own thoughts more clearly.
\end{itemize}

\subsection*{Conclusion: AI as a Tool -- But the Human Being Remains Responsible for Thinking}

Artificial intelligence is neither a replacement for
nor an opponent of human knowledge.
It is a \textbf{tool}\index{tool}
that can help with scientific communication --
\textbf{if it is used consciously, reflectively, and responsibly}.

In this respect as well,
this book understands itself as a contribution
to a new, enlightened way of dealing with technology in science.
Not because technology can do everything,
but because we must learn to use it meaningfully.

\vspace{1em}

\section*{Note on the Open Access Version}

This book is part of the Open Science initiative
``Christian \& Co-Pilot -- Math \& Physics.''

The complete, full-color PDF version is freely available at:

\begin{itemize}
	\item \href{https://mathandphysics.eu}{\texttt{https://mathandphysics.eu}}
	\item \href{https://zenodo.org/communities/christian-copilot}{\texttt{https://zenodo.org/communities/christian-copilot}}
\end{itemize}

The printed edition was created for comfortable reading,
long-term archiving,
and the development of an open-access library in print form.
By purchasing this book,
you support free scientific publication
and the idea of open, verifiable science.